\chapter*{Remerciements}
\addcontentsline{toc}{chapter}{Remerciements}

\vspace{0.5cm}

Je tiens a exprimer ma profonde gratitude a tous ceux qui ont contribue a la reussite de ce projet de fin d'etudes et a l'elaboration de ce rapport.

\vspace{0.5cm}

\textbf{Avant tout, je remercie Dieu le Tout-Puissant} de m'avoir donne la force, la sante et la volonte necessaires pour mener a bien ce projet. Sa benediction m'a accompagne tout au long de ce parcours.

\vspace{0.5cm}

\textbf{A mon encadrant au sein de CGI Fes,} pour son accueil chaleureux, sa confiance accordee dans le cadre de ce projet Zero Trust, ainsi que ses conseils pratiques et son accompagnement tout au long du stage. Son expertise du terrain et sa vision operationnelle ont ete essentielles pour adapter la solution aux contraintes reelles de l'entreprise.

\vspace{0.5cm}

\textbf{A mon encadrant academique,} pour ses conseils eclaires, son suivi rigoureux et sa disponibilite continue. Ses orientations methodologiques ont ete determinantes dans l'aboutissement de cette realisation.

\vspace{0.5cm}

\textbf{A l'equipe pedagogique} du departement Genie des Telecommunications et Reseaux (GTR), pour la qualite de la formation dispensee, qui m'a permis d'acquerir les competences necessaires a la conception et a l'implementation de cette architecture de cybersecurite.

\vspace{0.5cm}


\textbf{A mes collegues et camarades etudiants,} pour les echanges constructifs, l'entraide et le partage d'idees qui ont contribue a l'avancement de ce travail.

\vspace{0.5cm}

\textbf{A ma famille,} pour son soutien indefectible, sa patience et ses encouragements constants durant les nombreuses heures consacrees a ce projet.

\vspace{0.5cm}
Ce projet n'aurait pu aboutir sans cette convergence de competences, de conseils et d'encouragements. Il temoigne de l'importance de la collaboration dans le domaine de la cybersecurite, ou la mutualisation des connaissances est essentielle pour repondre aux defis securitaires actuels.

\vspace{1cm}

\newpage
